\begin{theorem}\label{night:P3K.45}
\StatementRecord{night:P3K.45}{false as stated; correction proved}{S6\_06\_BLIND: GPT-6 Astra (OpenAI, ChatGPT Pro mode, effort Pro), recorded re-derivation in the checkpoint dated 18 November 2026}{\cref{night:P3F.54}, \cref{night:P3F.56}, \cref{night:P3K.40}}\par\noindent\textbf{Physical-line correction, 17 November 2026.} The complete half-source conjugacy is retained. The reflected initial physical return is \(-1=27\), and the distinct fixed-edge flop is \(-7=21\). The old numerical labels are refuted by K68--69. The following earlier assertion and its argument are retained as the located historical claim; their corrected current consequences are proved in K68--69.\medskip\par 
Retain the selected-cusp source family and positive \(\rho\)-loop
of P3F.54/P3F.56, with fixed scalar, the actual III-star cap, and no
logarithmic modification on the coalescing finite pair. For every
allowed cusp fan \(T\), put \(J=\operatorname{diag}(-1,1)\)
on its marked two-dimensional lattice. The actual returns for
\(T\) and \(J(T)\) are conjugate by an orientation-preserving
diffeomorphism, through the literal half-source family map.
For the initial fan of P3F.56, \(J(T)\) is its nonfixed
diagonal-edge flop, and its positive return has oriented class
\[
                 [\Sigma_{\Phi,J(T)}]=1\in\mathbb Z/28.
\]
This gives neither a once-u endpoint on one fixed fan nor the
individual characteristic pair or integer derivative of the
reflected fan. The different fixed-edge flop has class23 by
P3K.40.
\end{theorem}
\begin{proof}
\begin{enumerate}[label=\arabic*.]
\item \par

\medskip\noindent\textbf{RF.1. The literal source, line and toric endpoint}

\par\nobreak

\item Retain the exact physical source equation and ordered node maps of
P3F.54. A positive half rho turn is one positive turn of the selected
u puncture. The square root in
\[
 \xi=\frac{y-X\sqrt{X+a(u)}}{\sqrt{X+u}},\qquad
 \eta=\frac{y+X\sqrt{X+a(u)}}{\sqrt{X+u}}
\]
changes sign, while the denominator root remains on its extending
branch. Thus the actual nodal endpoint exchanges \(\xi,\eta\).
It has the Tate source arrow \(Z\mapsto Z^{-1}\), with \(q_T\)
unchanged. The marked Picard point changes from \(P\) to
\(P^{-1}\).

\item We retain the line map. For the canonical theta ratio
\(f_P(Z)=\Theta(Z/P)/\Theta(Z)\), the exact identity in F54 is
\[
                   f_{P^{-1}}(Z^{-1})=P f_P(Z).
                                                               \tag{RF.1}
\]
The divisor-compatible arrow is consequently
\((Z,V)\mapsto(Z^{-1},PV)\). Reindex once by the new source
period \((Z',V')\mapsto(q_TZ',P^{-1}V')\). The two line factors
cancel, giving the actual annular endpoint
\[
                   (Z,V,q_T)\longmapsto(q_T/Z,V,q_T).
                                                               \tag{RF.2}
\]
Its square is literally the identity. The source translation and
its Hom map are intertwined by this same arrow: their source
multiplier changes to its inverse, while their scalar line factor
is retained. These are precisely the source and line transitions
of F54, not a replacement of the actual parameter family.

\item The III-star root charts, division section, unit root and finite
character extend over the full small u-disc. Their ordinary-to-cap
annuli continue with those coefficients. The half transport retains
these maps; its second iterate returns their literal charts. The
ordinary exterior is carried by the corresponding source and line
isomorphisms. Thus (RF.2) is compatible with every part of the actual
family, including its finite cap.

\item Here is the chart-ring map at the cusp. Write \(q_T=\kappa q\),
where \(q\) is the physical height character and \(\kappa\)
is the retained nonzero unit on the parameter annulus. For a target
character \(Z^mV^nq^\ell\), pullback by (RF.2) is
\[
          Z^mV^nq^\ell\longmapsto
             \kappa^m Z^{-m}V^nq^{\ell+m}.             \tag{RF.3}
\]
The integral character map is
\((m,n,\ell)\mapsto(-m,n,\ell+m)\). Its dual cone-height
map is
\[
                 (v_Z,v_V,h)\mapsto(-v_Z+h,v_V,h).
                                                               \tag{RF.4}
\]
On the height-one plane this is the affine reflection
\(p\mapsto Jp+e_1\). Periodicity identifies its translated fan
with \(J(T)\). Equation (RF.3) maps each target dual-cone
semigroup into the source semigroup with equality after the inverse
map. The extra factor \(\kappa^m\) is an actual holomorphic unit
on the full parameter annulus. Hence the monomial map and its
inverse extend over every actual unimodular chart. The period
actions are intertwined by the already displayed source and line
maps, so these chart maps descend to the compact models.

\item The physical delta action is multiplication in \(V\). It is
preserved literally by (RF.2), including its orientation. The
holomorphic chart maps preserve the complex orientation of the
six-dimensional fibres. No source-circle inversion or new line
character is substituted.

\par

\medskip\noindent\textbf{RF.2. The actual two-component family and its return arrows}

\par\nobreak

\item On the rho cover consider the two actual families with fans
\(T\) and \(J(T)\). The deck change \(\rho\mapsto-\rho\)
and (RF.2)--(RF.4) interchange them. Applying that arrow twice is
the identity on the source, line and toric charts, and the retained
finite-cap maps also return. Their disjoint union therefore descends
to a proper smooth family over the punctured u-disc whose fibre has
two connected components, labelled \(T,J(T)\) at the initial
point. This construction does not assert a once-u endpoint on either
one fixed fan.

\item Choose a smooth horizontal field on the disjoint cover family,
tangent to its prescribed collars, and average it under the actual
involution that exchanges the two families while taking rho to minus
rho. It descends to a horizontal field on the displayed u-family.
The physical delta action can be retained by an additional compact
circle average. Its transport along one positive u-loop has arrows
\[
           \alpha:X_T\longrightarrow X_{J(T)},\qquad
           \beta:X_{J(T)}\longrightarrow X_T.
                                                               \tag{RF.5}
\]
The domains and targets in (RF.5) are the actual components of that
same family. Restricting its second iterate to either component is
the physical return along one positive rho-loop in the corresponding
pulled-back family. Thus, with compatible representatives,
\[
                \Phi_T=\beta\alpha,\qquad
                \Phi_{J(T)}=\alpha\beta
                       =\alpha\Phi_T\alpha^{-1}.         \tag{RF.6}
\]
Any change of the horizontal field gives the usual isotopy through
its interpolated horizontal fields, with these same endpoint maps.
Equation (RF.6) therefore proves conjugacy of the actual positive
return classes. Since \(\alpha\) preserves the fibre orientation,
the oriented class in \(\Theta_7\) is the same for the two fans.
Explicitly, the already fixed identification of the oriented sphere
mapping-class group with \(\Theta_7=\mathbb Z/28\) is a group
homomorphism, and conjugation acts trivially in this abelian group.

\par

\medskip\noindent\textbf{RF.3. Identification of the particular reflected diagonal}

\par\nobreak

\item Use the exact two triangles of the initial physical-height atlas:
\[
       (0,e_1,e_2),\qquad(e_1,e_2,e_1+e_2),
                                                               \tag{RF.7}
\]
and all their lattice translates. Its three edge directions are
\(e_1,e_2,e_1-e_2\); the physical delta edge is \(e_2\).
The reflection \(J\) retains the two axis edge directions and
changes the third to \(e_1+e_2\), up to sign. It therefore replaces
exactly the diagonal \([e_1,e_2]\) of the unit parallelogram by
\([0,e_1+e_2]\). This is the periodic nonfixed diagonal-edge
flop, denoted \(C_2\) in the physical-height six-chart atlas.

\item P3F.56 proves the class of the original positive rho return is one.
Equations (RF.3)--(RF.7) consequently give the actual reflected-fan
return
\[
                    [\Sigma_{\Phi,J(T)}]=1
                             \quad\text{in }\mathbb Z/28.
                                                               \tag{RF.8}
\]
The individual derivative integer on \(C_2\) is still not computed
by this conjugacy. The separate fixed-edge flop of P3K.40 has class
23; its edge is \(C_0\), not the edge in (RF.7).

\item The certificate \path{code/part3/source_cusp_reflection.py} checks the
character/cone involution, physical delta retention, exact theta
line-factor cancellation, and the two labelled triangle sets.

\end{enumerate}
\end{proof}